\chapter{Quadratic Residues and Reciprocity}\label{ch:reciprocity}

\begin{center}\itshape
Is $219$ a square modulo the prime $383$? A direct search might square $191$ numbers.\\
With it, the answer (yes) takes four lines. Gauss called it the \emph{golden theorem} and returned to it again and again.
\end{center}


Square the nonzero residues modulo $7$. The results are $1,2,4$. Every other nonzero class fails to be a square. That short list is easy to check. The difficult question comes when the modulus is a large odd prime: must we try half its nonzero classes, or can we decide without finding the square root itself?

A remarkable connection compares two different primes. Modulo $7$, the number $11$ is the class $4$, hence a square. Modulo $11$, the number $7$ is \emph{not} a square: the nonzero squares there are $1,3,4,5,9$. The two answers disagree. Both primes are $3$ modulo $4$, and quadratic reciprocity explains why their signs must be opposite. If at least one prime is $1$ modulo $4$, the signs instead agree.

The theorem does not say that the primes are interchangeable. It gives a precise sign rule, with an exception controlled by their remainders modulo $4$. We will first encode the yes-or-no question using the Legendre symbol, then prove the rule by counting lattice points. After that we will calculate symbols without listing all possible squares.

One distinction remains important throughout: deciding that a square root exists modulo $p$ is not the same task as producing the root. A symbol answers the first question. A separate algorithm may be needed for the second.


\portrait{gauss}{Carl Friedrich Gauss (1777--1855). Portrait by Christian Albrecht Jensen, 1840. Public domain, via Wikimedia Commons.}

\begin{reminder}[Three questions, in the right order]
First: is a non-zero class a square? Euler's criterion answers by a modular
power. Second: how can we decide faster, or see the pattern as the prime
changes? Reciprocity replaces one prime by a smaller remainder. Third: if a
square root exists, how do we find it? That is a separate problem; a Legendre
symbol answers yes or no and does not hand over the root. We will keep these
questions apart, then connect them in a worked example.
\end{reminder}

\section{Quadratic residues and Euler's criterion}

% t1-pass2 motive 40c9ee6d
Start with the squares modulo $11$. Squaring $1,2,3,4,5$ gives $1,4,9,5,3$; the negatives give no new squares. So exactly five nonzero classes are squares. Euler's criterion turns that short list into a test that works for any odd prime $p$: one modular exponent tells us whether a nonzero class is a square. Before using it, understand why the nonzero squares come in pairs $x$ and $-x$ and why each pair has the same square.

Let $p$ be an odd prime. An integer $a$ with $p\nmid a$ is a \emph{quadratic residue} modulo $p$ if $x^2\equiv a\pmod p$ has a solution, and a \emph{non-residue} otherwise.

\begin{definition}
The \emph{Legendre symbol} is $\left(\frac ap\right)=1$ if $a$ is a quadratic residue, $-1$ if it is a non-residue, and $0$ if $p\mid a$.
\end{definition}

\begin{think}
List the squares modulo $11$: how many distinct non-zero classes do you get? Then compute $a^5\bmod11$ for $a=1,\dots,10$. Which values occur, and which $a$ give which value? Make a conjecture before reading the theorem.
\end{think}

\begin{theorem}[Euler's criterion]\label{thm:eulercrit}
$\left(\dfrac ap\right)\equiv a^{(p-1)/2}\pmod p$. Consequently the symbol is completely multiplicative in $a$, there are exactly $\frac{p-1}2$ residues and $\frac{p-1}2$ non-residues, and
\[
\left(\frac{-1}p\right)=(-1)^{(p-1)/2}:
\]
$-1$ is a square modulo $p$ if and only if $p\equiv1\pmod4$.
\end{theorem}
\begin{proof}
For $p\nmid a$, let $g$ be a primitive root (Chapter~\ref{ch:orders}) and write $a\equiv g^i\pmod p$. Then $a$ is a square if and only if $i$ is even (if $a=g^{2j}$, $x=g^j$ works; if $x=g^j$ then $x^2=g^{2j}$ and $2j\equiv i\pmod{p-1}$ forces $i$ even). Also $a^{(p-1)/2}=g^{i(p-1)/2}$, which is $1$ if $i$ is even and $g^{(p-1)/2}=-1$ if $i$ is odd (it squares to $1$ and is not $1$). Half of the exponents $0,\dots,p-2$ are even.
For $p\mid a$, both sides of the criterion are $0$, since $(p-1)/2\ge1$.
For any $a,b$, applying the criterion to $ab$ gives
$\left(\frac{ab}p\right)\equiv
\left(\frac ap\right)\left(\frac bp\right)\pmod p$.
If neither factor is zero these quantities belong to $\{1,-1\}$, whose
classes modulo the odd prime $p$ are distinct; if a factor is zero both
sides vanish. Thus the multiplicativity is an equality of integers, not just
a congruence. Finally apply the criterion with $a=-1$ to get the supplementary law.
\end{proof}

% t1-pass2 why
\paragraph{Why this works.}
The multiplicative group modulo $p$ is cyclic of even order, so squaring has a kernel of index $2$, the nonzero squares. The power $a^{(p-1)/2}$ is the unique character of order dividing $2$: Fermat says its square is $1$, and it equals $1$ on every square. Those two facts force it to be the Legendre symbol. The supplementary law for $-1$ is the same test at $a=-1$, where the power is $(-1)^{(p-1)/2}$.

\paragraph{Why exactly two roots?}
If $x^2\equiv y^2\pmod p$, then $p\mid(x-y)(x+y)$, so $x\equiv y$ or
$x\equiv-y$. Thus each non-zero square has exactly two roots, a pair
$\{x,-x\}$. There is no overlap inside such a pair because $p$ is odd and
$x\ne0$. The $p-1$ units therefore produce $(p-1)/2$ squares. This also
explains why zero is excluded from the word ``residue'' here: it has only one
root and the Legendre symbol treats it separately.

Modulo $11$, squaring $1,2,3,4,5$ gives $1,4,9,5,3$. The five non-zero
residues are therefore $1,3,4,5,9$. A compact experiment already displays the
criterion: $3^5\equiv1$, but $2^5=32\equiv-1$ modulo $11$.
Multiplying two non-residues gives a residue, since $(-1)(-1)=1$ at the level
of symbols. Addition has no such rule: $1$ and $3$ are residues but their sum
$4$ is too, while $1+1=2$ is a non-residue. The symbol is a multiplicative
sign, not an arithmetic fortune-teller.

\subsection*{A first character, and its cancellation}
Write $\chi(a)=\left(\frac ap\right)$. The identities just proved say that
$\chi$ is periodic modulo $p$, multiplicative, and zero precisely on multiples
of $p$. Such a function is called a \emph{Dirichlet character modulo $p$}.
Here we need only this particular real-valued example, not a general theory.

Its sum over a full residue system is zero. One proof counts the equally
numerous signs. Another exposes a useful mechanism: choose a non-residue $b$.
Multiplication by $b$ permutes the residues, so
\[
S:=\sum_{a=0}^{p-1}\chi(a)
=\sum_{a=0}^{p-1}\chi(ba)=\chi(b)S=-S,
\]
whence $S=0$. The cancellation is exact although the individual signs may
look irregular. It will let us calculate an expectation in
Exercise~\ref{ex:t1-residue-chance} without assuming independent signs.

% t1-pass2 examples 40c9ee6d
\begin{examplex}[residues modulo $11$]
The squares of $0,1,2,3,4,5$ modulo $11$ are $0,1,4,9,5,3$.
The nonzero residues are therefore $1,3,4,5,9$, and the non-residues are $2,6,7,8,10$.
Figure~\ref{fig:t1-residues} is that list as a strip. In particular $\bigl(\frac{2}{11}\bigr)=-1$, and $11\equiv 3\pmod 8$.
\end{examplex}

\begin{figure}[ht]
\centering
\fig{ch06-residues}
\caption{Quadratic residues modulo $11$, including $0$. Exactly five of the ten nonzero classes are squares, as $(11-1)/2=5$.}
\label{fig:t1-residues}
\end{figure}

\begin{examplex}[Euler's criterion for $p=13$]
Here $(p-1)/2=6$. Modulo $13$, $3^2=9$, $3^6=729$. Since $13\cdot 56=728$, one has $3^6\equiv 1$, so $3$ is a residue (indeed $4^2=16\equiv 3$). On the other hand $2^6=64\equiv 12\equiv -1$, so $2$ is a non-residue. The criterion sees the answer without listing every square.
\end{examplex}

\section{Gauss' lemma}

% t1-pass2 motive 92c448dd
A large modular power can answer a yes-or-no question, but it does not explain the answer very well. Gauss replaces the power by a count. Multiply $1,2,\ldots,(p-1)/2$ by a chosen residue $a$ and reduce each product to a signed representative between $-(p-1)/2$ and $(p-1)/2$. The parity of the negative representatives is the Legendre symbol. The proof is a carefully organised product in which every unwanted factor cancels legally.

\begin{theorem}[Gauss' lemma]\label{thm:gausslemma}
Let $p\nmid a$, and for $1\le k\le\frac{p-1}2$ let $r_k$ be the residue of $ka$ in $\left]-\frac p2,\frac p2\right[$. If $\mu$ is the number of negative $r_k$, then $\left(\frac ap\right)=(-1)^\mu$.
\end{theorem}
\begin{proof}
The $|r_k|$ lie in $\{1,\dots,\frac{p-1}2\}$ and are distinct: $ka\equiv\pm la$ implies $k\equiv\pm l$, so $k=l$ in this range. So they are a permutation of $1,\dots,\frac{p-1}2$. Multiplying all the $ka$:
$a^{(p-1)/2}\left(\frac{p-1}2\right)!\equiv\prod_kr_k=(-1)^\mu\left(\frac{p-1}2\right)!\pmod p$. Cancel the factorial and apply Euler's criterion.
\end{proof}

% t1-pass2 why
\paragraph{Why this works.}
Multiplying $1,\ldots,(p-1)/2$ by $a$ permutes a complete set of nonzero classes once the representatives are folded into $(-p/2,p/2)$. Each negative fold contributes one minus sign. Euler's criterion only needs that sign, so the symbol equals $(-1)^\mu$. The factorial cancels because it is not zero in the field $\Z/p\Z$.

\paragraph{A signed permutation, seen explicitly.}
Take $p=11$ and $a=3$. Reduce $3,6,9,12,15$ to the symmetric interval
$(-11/2,11/2)$; the representatives are
\[
 3,-5,-2,1,4.
\]
Their absolute values are a permutation of $1,2,3,4,5$, and there are two
minus signs. Hence $\left(\frac3{11}\right)=1$, agreeing with
$5^2\equiv3\pmod{11}$. For $a=2$, the representatives
$2,4,-5,-3,-1$ have three minus signs, so
$\left(\frac2{11}\right)=-1$.

The key cancellation in the proof is legal because none of
$1,\dots,(p-1)/2$ is divisible by $p$; the factorial is a unit.
To justify the distinctness carefully, $k\equiv-l\pmod p$ would require
$k+l=p$, impossible since $2\le k+l\le p-1$.
Gauss has turned a question about a square root into the parity of a count.
We count signs, not the size of the representatives.

\begin{corollary}\label{cor:two}
$\left(\dfrac2p\right)=(-1)^{(p^2-1)/8}$: $2$ is a square modulo $p$ if and only if $p\equiv\pm1\pmod8$.
\end{corollary}
\begin{proof}
For $a=2$, the numbers $2k$ with $1\le k\le\frac{p-1}2$ lie in $[2,p-1]$, and $r_k<0$ exactly when $2k>\frac p2$. So $\mu=\frac{p-1}2-\left\lfloor\frac p4\right\rfloor$. Checking $p=8m+1,3,5,7$ gives $\mu=2m,2m+1,2m+1,2m+2$, whose parity is that of $\frac{p^2-1}8$.
\end{proof}

% t1-pass2 why
\paragraph{Why this works.}
For $a=2$ the representatives that fall past $p/2$ are an interval of even numbers, and counting them is arithmetic, not reciprocity. The four odd classes modulo $8$ produce two even counts and two odd counts, which is the parity of $(p^2-1)/8$.

% t1-pass2 examples 92c448dd
\begin{examplex}[Gauss's count for $a=3$, $p=13$]
The products $3,6,9,12,15,18$ reduced into $(-13/2,13/2)$ are $3,6,-4,-1,2,5$.
Exactly two are negative, so $\bigl(\frac{3}{13}\bigr)=(-1)^2=1$.
And $6^2=36\equiv 10$, while $4^2=16\equiv 3$, which confirms the plus sign.
\end{examplex}

\begin{examplex}[the supplementary law at $p=17$ and $p=13$]
For $p=17\equiv 1\pmod 8$, $(p^2-1)/8=36$ is even, so $\bigl(\frac{2}{17}\bigr)=1$. Indeed $6^2=36\equiv 2$.
For $p=13\equiv 5\pmod 8$, $(169-1)/8=21$ is odd, so the symbol is $-1$, in agreement with $2^6\equiv -1$ above.
\end{examplex}

\section{The law of quadratic reciprocity}

% t1-pass2 motive ae6152e8
For distinct odd primes $p$ and $q$, compare the question whether $p$ is a square modulo $q$ with the question whether $q$ is a square modulo $p$. Their answers agree unless both primes are $3$ modulo $4$, in which case they disagree. This simple sentence hides the book's central surprise. The proof draws a rectangle of lattice points and counts it twice. Neither count alone gives reciprocity; matching them exposes the sign.

\begin{think}
Make a table: for the primes $p,q\in\{3,5,7,11,13\}$, $p\ne q$, record whether $p$ is a square modulo $q$. Now compare the entry for $(p,q)$ with the entry for $(q,p)$. When are they equal, and when do they differ? Look at $p$ and $q$ modulo $4$.
\end{think}

\begin{theorem}[quadratic reciprocity; Gauss, 1796]\label{thm:qr}
For distinct odd primes $p$ and $q$,
\[
\left(\frac pq\right)\left(\frac qp\right)=(-1)^{\frac{p-1}2\cdot\frac{q-1}2}.
\]
That is, $\left(\frac pq\right)=\left(\frac qp\right)$ unless $p\equiv q\equiv3\pmod4$, in which case $\left(\frac pq\right)=-\left(\frac qp\right)$.
\end{theorem}

\begin{figure}[ht]
\centering
\fig{ch06-reciprocity-table}
\caption{Each cell records $\left(\frac pq\right)$: row prime $p$, column
prime $q$. A plus means a non-zero square, a minus a non-square, and $0$
on the diagonal records divisibility. Compare cells reflected across the
diagonal: they disagree exactly when both primes are $3$ modulo $4$.
Signs carry the information even in black-and-white print.}
\label{fig:reciprocitytable}
\end{figure}

We give Eisenstein's proof, which counts lattice points. A second proof, with Gauss sums, is Exercise~\ref{ex:gausssum}.

\begin{lemma}[Eisenstein]\label{lem:eisenstein}
For odd $a$ with $p\nmid a$, $\left(\dfrac ap\right)=(-1)^{T(a,p)}$ where $T(a,p)=\displaystyle\sum_{k=1}^{(p-1)/2}\left\lfloor\frac{ka}p\right\rfloor$.
\end{lemma}
\begin{proof}
Write $ka=p\left\lfloor\frac{ka}p\right\rfloor+s_k$ with $0<s_k<p$. Let $A$ be the sum of the $s_k<\frac p2$ and $B$ the sum of the $\mu$ values $s_k>\frac p2$. In the proof of Gauss' lemma, the numbers $s_k$ ($s_k<\frac p2$) and $p-s_k$ ($s_k>\frac p2$) are $1,\dots,\frac{p-1}2$ in some order. Let $S=\sum_{k=1}^{(p-1)/2}k=\frac{p^2-1}8$. Then
\[
aS=p\,T(a,p)+A+B,\qquad S=A+\mu p-B.
\]
Subtracting, $(a-1)S=p\left(T(a,p)-\mu\right)+2B$. Modulo $2$, with $a$ and $p$ odd: $0\equiv T(a,p)-\mu$. Conclude with Gauss' lemma.
\end{proof}

% t1-pass2 why
\paragraph{Why this works.}
The floor sum $T(a,p)$ and Gauss's sign count $\mu$ differ by an even integer when $a$ is odd, because $(a-1)$ times the triangular number $(p^2-1)/8$ is even. Gauss's lemma already identifies the symbol with $(-1)^\mu$, hence with $(-1)^T$. The hypothesis that $a$ is odd is necessary: the same floors for $a=2$ can give the wrong sign.

\begin{figure}[ht]
\centering
\fig{ch06-eisenstein}
\caption{Eisenstein's count for $p=7$ and $q=11$. The $3\times5=15$ lattice points of the rectangle are split by the diagonal $y=\frac{11}7x$, which passes through none of them, into $T(11,7)=8$ and $T(7,11)=7$. So $\left(\frac{11}7\right)\left(\frac7{11}\right)=(-1)^{15}=-1$: indeed $11\equiv2^2$ is a square modulo $7$, but $7$ is not a square modulo $11$.}
\label{fig:eisenstein}
\end{figure}

\begin{think}
Before reading the proof, look at Figure~\ref{fig:eisenstein}. Count the points below the diagonal column by column, and those above it row by row. What do you get, and why does Eisenstein's lemma turn this picture into reciprocity?
\end{think}

\begin{proof}[Proof of Theorem~\ref{thm:qr}]
Count the points $(x,y)\in\Z^2$ with $1\le x\le\frac{p-1}2$ and $1\le y\le\frac{q-1}2$: there are $\frac{p-1}2\cdot\frac{q-1}2$. None lies on the line $qx=py$ (it would force $p\mid x$). Below the line ($y<\frac{qx}p$; note $\frac{qx}p<\frac q2$, so the bound $y\le\frac{q-1}2$ is automatic there) there are $\sum_x\left\lfloor\frac{qx}p\right\rfloor=T(q,p)$ points, and above it, symmetrically, $T(p,q)$. So $T(q,p)+T(p,q)=\frac{p-1}2\cdot\frac{q-1}2$, and Eisenstein's lemma gives the theorem.
\end{proof}

% t1-pass2 why
\paragraph{Why this works.}
The rectangle holds $(p-1)(q-1)/4$ lattice points and the diagonal misses every one of them, since a hit would be a multiple of both distinct primes. Points below the diagonal are $T(q,p)$ and points above are $T(p,q)$. Eisenstein's lemma turns those two integers into the two Legendre symbols, and their sum is the exponent of the reciprocity sign.

\paragraph{Why Eisenstein assumes that $a$ is odd.}
His proof ends by discarding $(a-1)S$ modulo $2$; this is even when $a$ is
odd. That assumption cannot simply be dropped. For $a=2,p=3$, the floor sum
is $\lfloor2/3\rfloor=0$, predicting a plus sign, whereas $2$ is not a square
modulo $3$. We handle factors of $2$ with the supplementary law, then apply
reciprocity only to odd numbers. Every hypothesis is doing a job.

\begin{example}
$383$ is prime and $219=3\cdot73$. Since $383\equiv3\pmod4$ and $3\equiv3\pmod4$: $\left(\frac3{383}\right)=-\left(\frac{383}3\right)=-\left(\frac23\right)=1$. Since $73\equiv1\pmod4$: $\left(\frac{73}{383}\right)=\left(\frac{383}{73}\right)=\left(\frac{18}{73}\right)=\left(\frac2{73}\right)\left(\frac9{73}\right)=1$ as $73\equiv1\pmod8$. So $\left(\frac{219}{383}\right)=1$.
\end{example}

\paragraph{A complete calculation, including the root.}
The previous example gives an existence decision. As a short alternative,
$18=2\cdot3^2$ modulo $73$, and $73\equiv1\pmod8$, so the only non-trivial
signs are already resolved by small remainders. What about producing a square
root of $219$ modulo $383$?

For a prime $p\equiv3\pmod4$ and a non-zero residue $a$, set
$r=a^{(p+1)/4}$. Then Euler's criterion gives
\[
 r^2=a^{(p+1)/2}=a\,a^{(p-1)/2}\equiv a\pmod p.
\]
Here $(383+1)/4=96$. Repeated squaring gives
\[
\begin{gathered}
219^2\equiv86,\quad219^4\equiv119,\quad219^8\equiv373,\\
219^{16}\equiv100,\quad219^{32}\equiv42,\quad219^{64}\equiv232
\pmod{383}.
\end{gathered}
\]
Thus $219^{96}\equiv232\cdot42\equiv169$ and the two roots are $169$ and
$383-169=214$. A final check is pleasantly short:
$169^2=28561=383\cdot74+219$. Existence, construction, verification:
each step has its own proof. The exponent formula is specific to
$p\equiv3\pmod4$; other primes require different methods.

% t1-pass2 examples ae6152e8
\begin{examplex}[a four-line decision]
Is $219$ a square modulo the prime $383$? Factor $219=3\cdot 73$.
Both $383$ and $3$ are $3$ modulo $4$, so $\bigl(\frac{3}{383}\bigr)=-\bigl(\frac{383}{3}\bigr)=-\bigl(\frac{2}{3}\bigr)=1$.
Since $73\equiv 1\pmod 4$, $\bigl(\frac{73}{383}\bigr)=\bigl(\frac{383}{73}\bigr)=\bigl(\frac{18}{73}\bigr)=\bigl(\frac{2}{73}\bigr)$.
And $73\equiv 1\pmod 8$, so the symbol is $1$. Thus $219$ is a square modulo $383$.
\end{examplex}

\begin{examplex}[the sign really flips]
Take $p=3$ and $q=7$, both $3$ modulo $4$. Then $\bigl(\frac{3}{7}\bigr)$ should be $-\bigl(\frac{7}{3}\bigr)$.
Squares modulo $7$ are $0,1,2,4$, so $3$ is a non-residue and the symbol is $-1$.
Modulo $3$, $7\equiv 1=1^2$, so $\bigl(\frac{7}{3}\bigr)=1$. The product of the symbols is $-1$, as the exponent $(3-1)/2\cdot(7-1)/2=3$ predicts.
\end{examplex}

\section{The Jacobi symbol}

% t1-pass2 motive d5918262
Reciprocity is most useful when it becomes a calculation. The Jacobi symbol extends the Legendre symbol's multiplication rules to odd composite denominators. This lets us reduce arguments by Euclid's algorithm without first factoring every large numerator. But the price is a new warning: Jacobi symbol $1$ does \emph{not} guarantee that a square root exists. We will compute an example where the test gives $1$ and the original congruence still has no solution.

For a positive odd integer $n=p_1\cdots p_r$ (primes with repetition), the \emph{Jacobi symbol} is $\left(\frac an\right)=\prod_i\left(\frac a{p_i}\right)$. It is multiplicative in $a$ (as each Legendre symbol is) and in $n$ (by definition). For $n=1$ the product is empty and the symbol is $1$, for every $a$.

\begin{proposition}\label{prop:jacobi}
For odd coprime $m,n>0$:
\[
\begin{gathered}\left(\frac{-1}n\right)=(-1)^{\frac{n-1}2},\qquad\left(\frac2n\right)=(-1)^{\frac{n^2-1}8},\qq\left(\frac mn\right)\left(\frac nm\right)=(-1)^{\frac{m-1}2\cdot\frac{n-1}2}.\end{gathered}
\]
\end{proposition}
\begin{proof}
For odd $a,b$: $\frac{ab-1}2-\frac{a-1}2-\frac{b-1}2=\frac{(a-1)(b-1)}2$ is even, and $\frac{a^2b^2-1}8-\frac{a^2-1}8-\frac{b^2-1}8=\frac{(a^2-1)(b^2-1)}8$ is even (each factor is divisible by $8$). So $u(n)=\frac{n-1}2$ and $v(n)=\frac{n^2-1}8$ are additive modulo $2$ over products of odd numbers. The first two formulas hold for primes (Theorem~\ref{thm:eulercrit}, Corollary~\ref{cor:two}), and both sides are multiplicative in $n$. For the third, with $m=\prod_iq_i$ and $n=\prod_jp_j$ (all distinct primes across the two lists, since $\gcd(m,n)=1$),
$\left(\frac mn\right)\left(\frac nm\right)=\prod_{i,j}\left(\frac{q_i}{p_j}\right)\left(\frac{p_j}{q_i}\right)=(-1)^{\sum_{i,j}u(q_i)u(p_j)}=(-1)^{u(m)u(n)}$, using Theorem~\ref{thm:qr} and additivity of $u$.
\end{proof}

% t1-pass2 why
\paragraph{Why this works.}
A Jacobi symbol is a product of Legendre symbols, so the prime law multiplies. The exponents $(n-1)/2$ and $(n^2-1)/8$ are additive modulo $2$ for odd factors, which is a one-line binomial check, and that is why the same closed formulas survive for composite denominators. What does not survive is the reading ``$1$ means a square exists''.

This gives a Euclid-like algorithm computing $\left(\frac an\right)$ \emph{without factoring}. Beware: $\left(\frac an\right)=1$ does not imply that $a$ is a square modulo a composite $n$ (e.g.\ $\left(\frac2{15}\right)=1$, but $2$ is not a square modulo $3$).

\paragraph{Computing a Jacobi symbol without factoring.}
Here is the algorithm behind the claim. Start with a positive odd denominator
$n$ and a sign $s=1$. At each stage reduce $a$ modulo $n$. If $a=0$, stop:
the answer is $0$ when $n>1$, and $s$ when $n=1$. If $a=1$, stop with
answer $s$. Otherwise remove every factor of $2$ from $a$, multiplying $s$
each time by $(-1)^{(n^2-1)/8}$. If this leaves $a=1$, stop with answer
$s$. The remaining $a>1$ is odd: swap $a$ and $n$, flipping $s$ when both
are $3$ modulo $4$, then repeat. After the reduction, the new denominator
is the old $a<n$, so the process terminates. The rule at $n=1$ is the
empty-product convention. The same reductions are Euclid's algorithm with
a few signs attached.
For non-coprime inputs, reciprocity still gives the required equality with
both symbols equal to $0$: an odd prime divides both numerator and denominator.
Removing factors of $2$ and taking Euclidean remainders preserve their gcd,
so this case ends at the zero rule with denominator greater than $1$.
This explains why the algorithm does not need a separate coprimality test.

For example, every swap in the following chain has positive sign because
at least one of its two odd numbers is $1$ modulo $4$:
\[
\begin{aligned}
\left(\frac{1001}{9907}\right)
 &=\left(\frac{898}{1001}\right)
 =\left(\frac{449}{1001}\right)\\
 &=\left(\frac{103}{449}\right)
 =\left(\frac{37}{103}\right)
 =\left(\frac{29}{37}\right)
 =\left(\frac8{29}\right)=-1.
\end{aligned}
\]
Here $\left(\frac2{1001}\right)=1$ because $1001\equiv1\pmod8$,
and $\left(\frac2{29}\right)=-1$ because $29\equiv5\pmod8$.
No factorisation of $9907$ was needed. A negative answer rules out a square;
a positive answer for a composite denominator need not certify one. The two
non-residues $2$ modulo $3$ and $2$ modulo $5$ give
$\left(\frac2{15}\right)=(-1)(-1)=1$, while $x^2\equiv2\pmod{15}$ has
no solution. A Jacobi symbol has forgotten some local information.

% t1-pass2 examples d5918262
\begin{examplex}[a composite denominator]
Compute $\bigl(\frac{2}{15}\bigr)=\bigl(\frac{2}{3}\bigr)\bigl(\frac{2}{5}\bigr)=(-1)\cdot(-1)=1$, since $3$ and $5$ are both $\pm 3$ modulo $8$.
Nevertheless $2$ is not a square modulo $15$: squares modulo $3$ are $0,1$, and $2\equiv -1$ is not among them. The Jacobi symbol has answered a different question.
\end{examplex}

\begin{examplex}[one Euclidean swap]
To compute $\bigl(\frac{19}{35}\bigr)$, reduce nothing yet: $19$ is already smaller than $35$, and both are $3$ modulo $4$, so the symbol equals $-\bigl(\frac{35}{19}\bigr)=-\bigl(\frac{16}{19}\bigr)$.
But $16=4^2$, so $\bigl(\frac{16}{19}\bigr)=1$ and the original symbol is $-1$.
No prime factor of $35$ was required.
\end{examplex}

\begin{beauty}
Where this leads. The Legendre symbol is the Dirichlet character Volume~6 uses for the primes in an arithmetic progression, and the supplementary laws are the characters modulo $4$ and $8$. Quadratic reciprocity is the abelian reciprocity law in its smallest case; Volume~13 rebuilds it as the splitting law in cyclotomic fields and then as class field theory. The algorithm that computes a Jacobi symbol without factoring is also the algorithm that evaluates those characters at a prime.
\end{beauty}

\subsection*{Proof workshop: let the small primes test the sign}

For $p=7$, the distinct nonzero squares are
$1^2\equiv1$, $2^2\equiv4$ and $3^2\equiv2$.
Thus $\left(\frac{11}{7}\right)=\left(\frac47\right)=1$.

For $q=11$, the squares of $1,2,3,4,5$ are
$1,4,9,5,3$ modulo $11$. The residue $7$ is absent,
so $\left(\frac7{11}\right)=-1$. Their product is $-1$.

The reciprocity sign is
\[
(-1)^{((7-1)/2)((11-1)/2)}=(-1)^{3\cdot5}=-1.
\]
The calculation agrees with the theorem. It does not
prove the theorem: we tested a single pair of primes.
The proof must explain that sign for \emph{every} pair
of distinct odd primes. Lattice-point counting supplies
the missing general argument.


\section{Exercises}

\begin{exercise}[$\star$]\label{ex:minus3}
Show that for a prime $p>3$, $-3$ is a square modulo $p$ if and only if $p\equiv1\pmod3$.
\end{exercise}

\begin{exercise}[$\star$]\label{ex:onemod4}
Show that every odd prime factor of $n^2+1$ is $\equiv1\pmod4$, and deduce that there are infinitely many primes $p\equiv1\pmod4$.
\end{exercise}

\begin{exercise}[$\star\star$]\label{ex:minus2}
For an odd prime $p$, show that $\left(\frac{-2}p\right)=1$ if and only if $p\equiv1$ or $3\pmod8$. Deduce that a prime of the form $x^2+2y^2$ is $2$ or $\equiv1,3\pmod8$.
\end{exercise}

\begin{exercise}[$\star\star$]\label{ex:mersfactor}
Let $p\equiv3\pmod4$ be prime with $q=2p+1$ also prime. Show that $q\mid2^p-1$. Deduce that $2^{11}-1$, $2^{23}-1$ and $2^{251}-1$ are composite.
\end{exercise}

\begin{exercise}[$\star\star$]\label{ex:consecutive}
Let $p$ be an odd prime. Show that $\displaystyle\sum_{a=1}^{p-1}\left(\frac{a(a+1)}p\right)=-1$. Deduce the number of $a\in\{1,\dots,p-2\}$ such that $a$ and $a+1$ are both quadratic residues.
\end{exercise}

\begin{exercise}[$\star\star$]\label{ex:pepin}
\emph{P\'epin's test.} Let $\mathcal F_n=2^{2^n}+1$ with $n\ge1$. Show that $\mathcal F_n$ is prime if and only if $3^{(\mathcal F_n-1)/2}\equiv-1\pmod{\mathcal F_n}$.
\end{exercise}

\begin{exercise}[$\star\star\star$]\label{ex:gausssum}
\emph{Reciprocity through Gauss sums.} Let $\zeta=e^{2i\pi/p}$, $p^*=\left(\frac{-1}p\right)p$, and $g=\sum_{a=1}^{p-1}\left(\frac ap\right)\zeta^a$.
\begin{enumerate}[label=(\alph*)]
\item Show that $g^2=p^*$. \emph{Hint: substitute $b=ac$ in $g^2=\sum_{a,b}\left(\frac{ab}p\right)\zeta^{a+b}$.}
\item Let $q\ne p$ be an odd prime. In the ring $\Z[\zeta]$, show that $g^q\equiv\left(\frac qp\right)g\pmod{q\Z[\zeta]}$.
\item Show that $q\Z[\zeta]\cap\Z=q\Z$, using that $\Phi_p$ is irreducible over $\Q$ (Volume~1, the polynomial exercise on prime cyclotomic irreducibility). Deduce $\left(\frac{p^*}q\right)=\left(\frac qp\right)$, and then Theorem~\ref{thm:qr}.
\item Show that $g=\sum_{x=0}^{p-1}\zeta^{x^2}$, the sum $G$ of Volume~1, the complex-number exercise on quadratic Gauss sums. What does (a) say about $G$?
\end{enumerate}
\end{exercise}

\begin{exercise}[$\star\star$]\label{ex:t1-residue-chance}
\emph{Residue signs and chance.} Choose $U$ uniformly from
$\{0,\dots,p-1\}$ for an odd prime $p$, and put
$X=\left(\frac Up\right)$, $Y=\left(\frac{U+1}p\right)$.
Find the distributions, expectations and variances of $X,Y$.
Using Exercise~\ref{ex:consecutive}, prove
$\operatorname{Cov}(X,Y)=-1/p$. Are the two signs independent?
\end{exercise}

\begin{exercise}[$\star$]\label{ex:t1-res11}
List the quadratic residues modulo $11$. Decide $\bigl(\frac{6}{13}\bigr)$ and $\bigl(\frac{6}{17}\bigr)$ by the supplementary laws and reciprocity, and check one of them by listing squares.
\end{exercise}

\begin{beauty}
Quadratic reciprocity links two questions that have no visible connection: whether $p$ is a square modulo $q$, and whether $q$ is a square modulo $p$. Many proofs are known, each exposing a different structure. Its generalisations (cubic and biquadratic reciprocity, which live in the rings of Chapter~\ref{ch:gaussian}, Artin reciprocity, and the Langlands program) have driven number theory for two centuries. It also powers the Solovay--Strassen primality test and the computation of square roots modulo $p$.
\end{beauty}

\begin{exercise}[$\star\star$]\label{pb:primsum}
Let $p$ be an odd prime.
\begin{enumerate}[label=(\alph*)]
\item Show that the sum of the primitive roots modulo $p$ in $\{1,\dots,p-1\}$ is congruent to $\mu(p-1)$ modulo $p$.
\item Show that their product is congruent to $1$ modulo $p$ when $p>3$.
\item More generally, show that the sum of the $k$-th powers of the primitive roots is congruent modulo $p$ to the
Ramanujan sum $c_{p-1}(k)$ of Problem~\ref{pb:ramsum}.
\end{enumerate}
\end{exercise}

\begin{exercise}[$\star\star$]\label{pb:phimodp}
\begin{enumerate}[label=(\alph*)]
\item Let $p$ be a prime not dividing $n$, and $f=\operatorname{ord}_n(p)$. Show that $\Phi_n$ modulo $p$ is a product of
$\varphi(n)/f$ distinct irreducible polynomials of $\mathbb F_p[X]$, all of degree $f$.
\item Deduce that $X^4+1$ is irreducible over $\Q$ but factors modulo every prime.
\item Let $n=p^am$ with $a\ge1$ and $p\nmid m$. Show that $\Phi_n\equiv\Phi_m^{\varphi(p^a)}\pmod p$.
\end{enumerate}
\end{exercise}

\begin{exercise}[$\star$]\label{pb:fermatnum}
Let $F_n=2^{2^n}+1$ be the Fermat numbers.
\begin{enumerate}[label=(\alph*)]
\item Show that every prime factor of $F_n$, $n\ge2$, is congruent to $1$ modulo $2^{n+2}$.
\item Show that $641$ divides $F_5$, without dividing, using $641=5\cdot2^7+1=5^4+2^4$.
\item Show that two distinct Fermat numbers are coprime, and deduce that there are infinitely many primes.
\end{enumerate}
\end{exercise}

\begin{exercise}[$\star\star$]\label{pb:carmichael}
A composite $n$ is a \emph{Carmichael number} if $a^{n-1}\equiv1\pmod n$ for every $a$ prime to $n$.
\begin{enumerate}[label=(\alph*)]
\item Show that a Carmichael number is odd, squarefree, and that $p-1\mid n-1$ for every prime $p\mid n$.
\item Show that a Carmichael number has at least three prime factors.
\item Show that if $6k+1$, $12k+1$ and $18k+1$ are all prime, their product is a Carmichael number. Find the smallest
example.
\end{enumerate}
\end{exercise}

\begin{exercise}[$\star\star$]\label{pb:pratt}
\begin{enumerate}[label=(\alph*)]
\item Show that $n\ge3$ is prime if and only if there is an $a$ with $a^{n-1}\equiv1\pmod n$ and $a^{(n-1)/q}\not\equiv1\pmod n$
for every prime $q\mid n-1$.
\item A \emph{certificate} for an odd prime $p$ is such an $a$, the factorisation of $p-1$, and a certificate for each odd
prime factor of $p-1$, and so on down. Let $f(p)$ be the number of primes in the whole tree, with $f(2)=1$ and
$f(p)=1+\sum_{q\mid p-1}f(q)$, the sum over the distinct primes $q$. Show that $f(p)\le2\log_2p-1$ for every odd prime $p$.
\item Deduce that a prime of $k$ binary digits has a certificate that can be checked with a number of operations on
$k$-digit numbers bounded by a polynomial in $k$.
\end{enumerate}
\end{exercise}

\begin{exercise}[Journey]\label{pb:rabin}
\emph{The Miller--Rabin test lies at most one time in four.} Let $n$ be an odd composite number, $n-1=2^sd$ with $d$ odd.
An $a$ prime to $n$ is a \emph{strong liar} if $a^d\equiv1\pmod n$ or $a^{2^jd}\equiv-1\pmod n$ for some $0\le j<s$. Write
$n=\prod_{i=1}^rp_i^{e_i}$ and $p_i-1=2^{\nu_i}m_i$ with $m_i$ odd, and let $\nu=\min_i\nu_i$.
\begin{enumerate}[label=(\alph*)]
\item In a cyclic group $C$ of even order $m$, with $-1$ its element of order $2$, show that $x^t=1$ has $\gcd(t,m)$
solutions and $x^t=-1$ has $\gcd(2t,m)-\gcd(t,m)$ solutions.
\item (Monier) Show that the number of strong liars is
\[
\prod_{i=1}^r\gcd(d,m_i)\cdot\Bigl(1+\frac{2^{r\nu}-1}{2^r-1}\Bigr).
\]
\item Deduce that at most $\frac14\varphi(n)$ of the units are strong liars, unless $n=9$. Check that $15$ attains $\frac14$.
\end{enumerate}
\end{exercise}

\begin{exercise}[$\star\star\star$]\label{pb:waring}
\emph{Two $k$-th powers.} Let $p$ be prime, $k\ge2$ a divisor of $p-1$, and $N(c)$ the number of $(x,y)\in\mathbb F_p^2$ with
$x^k+y^k=c$.
\begin{enumerate}[label=(\alph*)]
\item Show that for $c\ne0$,
\[
|N(c)-p|\le(k-1)(k-2)\sqrt p+k-1.
\]
\item Deduce that every residue modulo $p$ is a sum of two $k$-th powers when $p\ge k^4$.
\item What happens when $k$ does not divide $p-1$?
\end{enumerate}
\end{exercise}

\begin{exercise}[$\star\star$]\label{pb:squaresmodpk}
\emph{Squares modulo $p^k$.} Let $S(m)$ be the number of squares in $\Z/m\Z$ ($0$ included).
\begin{enumerate}[label=(\alph*)]
\item For an odd prime $p$, show that $S(p^k)=1+\sum_{0\le j<k/2}\frac{p-1}2p^{k-2j-1}$, and that $S(p^k)/p^k\to\frac p{2(p+1)}$.
\item Show that $S(2^k)/2^k\to\frac16$.
\item Show that $S$ is multiplicative, and deduce that the proportion of squares modulo $m$ can be made as small as we like.
What does this say about an integer that is a square modulo every $m$?
\end{enumerate}
\end{exercise}

\begin{figure}[ht]
\centering
\fig{ch03-liars}
\caption{For each odd composite $n<1500$, the proportion of units that are Euler liars (coral) and strong liars (charcoal). The
coral dots stay below $\frac12$ (Problem~\ref{pb:euler}), the charcoal ones below $\frac14$ (Problem~\ref{pb:rabin}), except for
$n=9$, at $\frac13$. Most composite numbers have almost no liars at all: the bounds are worst cases.}
\label{fig:liars}
\end{figure}

\section{Hints}

\begin{hint}{ex:minus3}
$\left(\frac{-3}p\right)=\left(\frac{-1}p\right)\left(\frac3p\right)$. Apply Euler's criterion to the first factor and reciprocity to the second: the signs cancel.
\end{hint}
\begin{hint}{ex:onemod4}
$p\mid n^2+1$ means that $-1$ is a square modulo $p$. For infinitude, take $n=2p_1\cdots p_k$.
\end{hint}
\begin{hint}{ex:minus2}
Multiply the two supplementary laws and tabulate $p\bmod8$. For the second part, if $p=x^2+2y^2$ then $p\nmid y$, and $(xy^{-1})^2\equiv-2$.
\end{hint}
\begin{hint}{ex:mersfactor}
$q\mid2^p-1$ means $2^{(q-1)/2}\equiv1\pmod q$, that is, $2$ is a square modulo $q$. What is $q\bmod8$?
\end{hint}
\begin{hint}{ex:consecutive}
For $a\ne0$, $\left(\frac{a(a+1)}p\right)=\left(\frac{1+a^{-1}}p\right)$. Which values does $1+a^{-1}$ take as $a$ runs over the units? For the count, express the indicator of ``$a$ and $a+1$ are residues'' with $\frac{1+\chi(a)}2\cdot\frac{1+\chi(a+1)}2$.
\end{hint}
\begin{hint}{ex:pepin}
If $\mathcal F_n$ is prime, compute $\left(\frac3{\mathcal F_n}\right)$ by reciprocity, using $\mathcal F_n\bmod3$ and $\mathcal F_n\bmod4$. For the converse, the congruence gives the order of $3$ modulo any prime factor $p$, and it must divide $p-1$.
\end{hint}
\begin{hint}{ex:gausssum}
(b) In characteristic $q$, $(x+y)^q\equiv x^q+y^q$, so $g^q\equiv\sum_a\left(\frac ap\right)\zeta^{aq}$; substitute $b=aq$. (c) Multiply (b) by $g$ and use (a); then both sides are integers. Do not cancel $g$ inside the quotient ring before proving that it is invertible there. Instead descend to an integer congruence and cancel $p^*$ modulo $q$. (d) Count the solutions of $x^2\equiv a$ with $1+\left(\frac ap\right)$.
\end{hint}

\begin{hint}{ex:t1-residue-chance}
Zero occurs once and each nonzero sign occurs $(p-1)/2$ times.
Sum $\chi(a)\chi(a+1)$ over all $a$; the missing term at $a=0$ is zero.
\end{hint}

\begin{hint}{ex:t1-res11}
Squares modulo $11$ need only the representatives $0,\ldots,5$. Factor $6=2\cdot 3$, and remember that the sign flips only when both primes are $3$ modulo $4$.
\end{hint}

\begin{hint}{pb:primsum}
In $\mathbb F_p$, the primitive roots are the $\gamma^a$ with $\gcd(a,p-1)=1$, for one primitive root $\gamma$. Copy the
proof of Problem~\ref{pb:ramsum}(a) with $\gamma$ in place of $e^{2i\pi/q}$.
\end{hint}

\begin{hint}{pb:phimodp}
(a) $X^n-1$ has distinct roots modulo $p$, so the roots of $\Phi_n$ modulo $p$ have order exactly $n$. An element of order $n$
lies in $\mathbb F_{p^f}$ exactly when $n\mid p^f-1$. (c) Use $\Phi_m(x^p)\equiv\Phi_m(x)^p\pmod p$.
\end{hint}

\begin{hint}{pb:fermatnum}
(a) The order of $2$ modulo $p$ is $2^{n+1}$, and $2$ is a square modulo $p$ when $p\equiv1\pmod8$. (c) $F_0F_1\cdots F_{n-1}=F_n-2$.
\end{hint}

\begin{hint}{pb:carmichael}
(a) Try $a=-1$; then use a primitive root modulo $p^2$ or modulo $p$, glued by the Chinese remainder theorem. (b) If $n=pq$
with $p<q$, compute $n-1$ modulo $q-1$.
\end{hint}

\begin{hint}{pb:pratt}
(a) An element of order $n-1$ forces $\varphi(n)=n-1$. (b) Induction, using $\prod q_i\le\frac{p-1}2$ over the odd primes $q_i\mid p-1$.
\end{hint}

\begin{hint}{pb:rabin}
(b) Use the Chinese remainder theorem and the cyclic groups $(\Z/p_i^{e_i}\Z)^*$; note that $p_i\nmid d$. (c) Compare with
$\varphi(n)=\prod p_i^{e_i-1}2^{\nu_i}m_i$. For $r=2$, the hardest case is $n=p_1p_2$ with $\nu_1=\nu_2$.
\end{hint}

\begin{hint}{pb:waring}
Write the number of $k$-th roots of $u$ as a sum of characters, and use Jacobi sums of absolute value $\sqrt p$.
\end{hint}

\begin{hint}{pb:squaresmodpk}
Write a non-zero square as $p^{2j}u$ with $u$ a unit; the units that are squares form a subgroup of index $2$ ($4$ for $p=2$).
\end{hint}

\section{Solutions}

\begin{solution}{ex:minus3}
\idea Split $-3=(-1)\cdot3$ by multiplicativity: $\left(\frac{-1}p\right)$ is given by Euler's criterion and $\left(\frac3p\right)$ by reciprocity. The two signs $(-1)^{(p-1)/2}$ cancel, leaving $\left(\frac p3\right)$, which only depends on $p\bmod3$.
\writeup $\left(\frac{-3}p\right)=\left(\frac{-1}p\right)\left(\frac3p\right)=(-1)^{\frac{p-1}2}\cdot(-1)^{\frac{p-1}2}\left(\frac p3\right)=\left(\frac p3\right)$, by reciprocity with the exponent $\frac{p-1}2\cdot\frac{3-1}2=\frac{p-1}2$. And $\left(\frac p3\right)=1$ iff $p\equiv1\pmod3$ (the squares modulo $3$ are $0,1$).
\end{solution}

\begin{solution}{ex:onemod4}
\idea $p\mid n^2+1$ says exactly that $-1$ is a square modulo $p$, which Euler's criterion translates into $p\equiv1\pmod4$. For infinitude, imitate Euclid with $n^2+1$ in place of $n+1$, choosing $n$ divisible by all the known primes $\equiv1\pmod4$ (and by $2$, to make $n^2+1$ odd).
\writeup If $p\mid n^2+1$ with $p$ odd, $-1\equiv n^2$ is a square modulo $p$, so $p\equiv1\pmod4$ by Euler's criterion. If $p_1,\dots,p_k$ were all the primes $\equiv1\pmod4$, then $N=(2p_1\cdots p_k)^2+1>1$ is odd, and its prime factors are $\equiv1\pmod4$ but differ from all the $p_i$ (which would divide $1$).
\end{solution}

\begin{solution}{ex:minus2}
\idea Multiply the two supplementary laws, $\left(\frac{-1}p\right)$ and $\left(\frac2p\right)$, and tabulate modulo $8$. For the second part: $p=x^2+2y^2$ gives $x^2\equiv-2y^2$, so $-2$ is a square as soon as $y$ is invertible.
\writeup $\left(\frac{-2}p\right)=(-1)^{\frac{p-1}2}(-1)^{\frac{p^2-1}8}$. For $p\equiv1,3,5,7\pmod8$ the two signs are $(+,+),(-,-),(+,-),(-,+)$, so the product is $1$ exactly for $p\equiv1,3$. If $p=x^2+2y^2$ is odd, then $p\nmid y$ (else $p\mid x$ and $p^2\mid p$), and $\left(xy^{-1}\right)^2\equiv-2$.
\end{solution}

\begin{solution}{ex:mersfactor}
\idea $q\mid2^p-1$ means $2^{(q-1)/2}\equiv1\pmod q$, which by Euler's criterion means $2$ is a square modulo $q$. And $\left(\frac2q\right)$ only depends on $q\bmod8$: here $q=2p+1$ with $p\equiv3\pmod4$ gives $q\equiv7\pmod8$.
\writeup $q=2p+1\equiv7\pmod8$, so $\left(\frac2q\right)=1$ (Corollary~\ref{cor:two}) and Euler's criterion gives $2^{(q-1)/2}=2^p\equiv1\pmod q$. For $p=11,23,251$ the numbers $q=23,47,503$ are prime, and $q<2^p-1$, so $q$ is a proper factor.
\end{solution}

\begin{solution}{ex:consecutive}
\idea Sums of Legendre symbols of polynomials become tractable after a substitution that turns the polynomial into something linear. Factor out the square $a^2$: $\chi(a(a+1))=\chi(1+a^{-1})$, and $1+a^{-1}$ runs over almost all residues. For the count, the indicator of ``$a$ is a residue'' is $\frac{1+\chi(a)}2$, so the count of pairs is a sum of products of such indicators, which expands into sums we know.
\writeup For $1\le a\le p-1$, $\left(\frac{a(a+1)}p\right)=\left(\frac{a^2(1+a^{-1})}p\right)=\left(\frac{1+a^{-1}}p\right)$. As $a$ runs over the non-zero residues, $1+a^{-1}$ runs over all residues except $1$. Since $\sum_{c=0}^{p-1}\left(\frac cp\right)=0$ (equally many residues and non-residues), the sum is $0-\left(\frac1p\right)=-1$.
Let $N$ be the number of $a\in[1,p-2]$ with both $a$, $a+1$ residues. With $\chi=\left(\frac\cdot p\right)$, $\frac14(1+\chi(a))(1+\chi(a+1))$ is $1$ for such $a$ and $0$ for other $a\in[1,p-2]$. Summing over $a\in[1,p-2]$: $\sum\chi(a)=-\chi(p-1)=-\chi(-1)$, $\sum\chi(a+1)=-\chi(1)=-1$, and $\sum\chi(a)\chi(a+1)=-1-\chi(p-1)\chi(p)=-1$. So $N=\frac14\left(p-2-\chi(-1)-1-1\right)=\frac{p-4-\left(\frac{-1}p\right)}4$: $\frac{p-5}4$ if $p\equiv1\pmod4$ and $\frac{p-3}4$ if $p\equiv3\pmod4$.
\end{solution}

\begin{solution}{ex:pepin}
\idea If $\mathcal F_n$ is prime, Euler's criterion says $3^{(\mathcal F_n-1)/2}\equiv\left(\frac3{\mathcal F_n}\right)$, and reciprocity computes this symbol from $\mathcal F_n\bmod3$. For the converse, the congruence gives the exact order of $3$ (a power of $2$ whose half fails), and an element of order $\mathcal F_n-1$ forces $\varphi(\mathcal F_n)=\mathcal F_n-1$, i.e.\ primality. This is the template of Lucas-type primality proofs.
\writeup If $\mathcal F_n$ is prime: $\mathcal F_n\equiv1\pmod4$ and $\mathcal F_n\equiv2\pmod3$ (as $2^{2^n}\equiv1\pmod 3$, $2^n$ being even), so by reciprocity $\left(\frac3{\mathcal F_n}\right)=\left(\frac{\mathcal F_n}3\right)=\left(\frac23\right)=-1$, and Euler's criterion gives the congruence. Conversely, if $3^{(\mathcal F_n-1)/2}\equiv-1$, then $3^{\mathcal F_n-1}\equiv1$. This already forces $\gcd(3,\mathcal F_n)=1$, so $\operatorname{ord}(3)$ divides $\mathcal F_n-1=2^{2^n}$ but not its half: it equals $\mathcal F_n-1$. Since the order divides $\varphi(\mathcal F_n)\le\mathcal F_n-1$, we get $\varphi(\mathcal F_n)=\mathcal F_n-1$, so every $1\le k<\mathcal F_n$ is prime to $\mathcal F_n$, and $\mathcal F_n$ is prime.
\end{solution}

\begin{solution}{ex:gausssum}
\idea The Gauss sum $g$ is an algebraic number whose square is $\pm p$, so it ``is'' $\sqrt{p^*}$ inside $\Z[\zeta]$. Then $\left(\frac{p^*}q\right)$, which by Euler's criterion is $(p^*)^{(q-1)/2}=g^{q-1}$ modulo $q$, can be computed in a second way: raising to the power $q$ modulo $q$ is additive (the Frobenius), and it permutes the $\zeta^a$, producing $\left(\frac qp\right)$. Comparing the two computations gives reciprocity. The only subtle point is to come back from congruences in $\Z[\zeta]$ to congruences in $\Z$, which needs the linear independence of the powers of $\zeta$.
\writeup (a) With $\chi=\left(\frac\cdot p\right)$,
\[
g^2=\sum_{a,b\ne0}\chi(ab)\zeta^{a+b}.
\] Put $b=ac$ ($c$ runs over non-zero residues for fixed $a$): $\chi(ab)=\chi(a^2c)=\chi(c)$, so
\[g^2=\sum_{c\ne0}\chi(c)\sum_{a\ne0}\zeta^{a(1+c)}.\]
The inner sum is $p-1$ if $c\equiv-1$ and $-1$ otherwise (Volume~1, the roots-of-unity orthogonality proposition). Hence
\[\begin{aligned}g^2&=\chi(-1)(p-1)-\sum_{c\ne0,-1}\chi(c)\eqbr=\chi(-1)(p-1)+\chi(-1)=\chi(-1)p=p^*,\end{aligned}\]
using $\sum_{c\ne0}\chi(c)=0$.

(b) In any commutative ring, $(x+y)^q\equiv x^q+y^q$ modulo $q$ (the binomial coefficients $\binom qj$, $0<j<q$, are multiples of $q$), and by induction for sums of several terms. So $g^q\equiv\sum_a\chi(a)^q\zeta^{aq}=\sum_a\chi(a)\zeta^{aq}$, as $\chi(a)^q=\chi(a)$ ($q$ odd). Substituting $b=aq$: $\chi(a)=\chi(bq^{-1})=\chi(q)\chi(b)$ (as $\chi(q^{-1})=\chi(q)$), so the sum is $\chi(q)g$.

(c) First, $q\Z[\zeta]\cap\Z=q\Z$. The numbers $1,\zeta,\dots,\zeta^{p-2}$ are linearly independent over $\Q$: a relation would give a non-zero $R\in\Q[X]$ of degree $\le p-2$ with $R(\zeta)=0$; then $\gcd(R,\Phi_p)$, which is a $\Q[X]$-combination of $R$ and $\Phi_p$ (B\'ezout in $\Q[X]$), also vanishes at $\zeta$, so it is a non-constant proper divisor of $\Phi_p$, contradicting the irreducibility of $\Phi_p$ (Volume~1, the polynomial exercise on prime cyclotomic irreducibility). Every element of $\Z[\zeta]$ is $\sum_{i\le p-2}c_i\zeta^i$ with $c_i\in\Z$ (reduce higher powers with $\zeta^{p-1}=-1-\zeta-\dots-\zeta^{p-2}$). If an integer $n$ equals $q\sum c_i\zeta^i$, uniqueness of coordinates gives $n=qc_0$.

Now multiply (b) by $g$: $g^{q+1}\equiv\chi(q)g^2$, that is $(p^*)^{(q+1)/2}\equiv\chi(q)p^*$ modulo $q\Z[\zeta]$. Both sides are integers, so the congruence holds modulo $q$ in $\Z$. As $p^*$ is invertible modulo $q$: $(p^*)^{(q-1)/2}\equiv\chi(q)\pmod q$, and by Euler's criterion the left side is $\left(\frac{p^*}q\right)$; both are $\pm1$ and $q>2$, so they are equal. Finally $\left(\frac{p^*}q\right)=\left(\frac{-1}q\right)^{\frac{p-1}2}\left(\frac pq\right)=(-1)^{\frac{q-1}2\cdot\frac{p-1}2}\left(\frac pq\right)$, which is Theorem~\ref{thm:qr}.

(d) The number of $x$ with $x^2\equiv a$ is $1+\chi(a)$ (two roots for a residue, none for a non-residue, one for $0$), so $\sum_x\zeta^{x^2}=\sum_a(1+\chi(a))\zeta^a=\sum_a\zeta^a+g=g$. So $G^2=p^*$: $G=\pm\sqrt p$ if $p\equiv1\pmod4$ and $\pm i\sqrt p$ if $p\equiv3\pmod4$, consistent with $|G|^2=p$. A numerical calculation suggests that the sign is always $+$; this is a much harder theorem of Gauss, proved in Volume~5 with Fourier analysis.
\end{solution}

\begin{solution}{ex:t1-residue-chance}
\idea Uniform sampling converts an exact finite character sum into an
expectation. Marginal balance does not imply independence: examine the
product, rather than guessing from a picture of the signs.
\writeup Each variable is $0$ with probability $1/p$, and each of $1,-1$
with probability $(p-1)/(2p)$. Translation by $1$ permutes the sample space,
so the distributions agree. Thus $\mathbb E X=\mathbb E Y=0$ and
$\operatorname{Var}(X)=\operatorname{Var}(Y)=(p-1)/p$.
Exercise~\ref{ex:consecutive} gives
\[
\mathbb E(XY)=\frac1p\sum_{a=0}^{p-1}\chi(a(a+1))=-\frac1p.
\]
Consequently the covariance is $-1/p$. Independence would imply
$\mathbb E(XY)=(\mathbb E X)(\mathbb E Y)=0$, contradicting this value. Even a balanced
multiplicative sign pattern need not give independent neighbouring signs.
\end{solution}

\begin{solution}{ex:t1-res11}
\idea The list is five squares. Each symbol factors, and only the factor $3$ can produce a reciprocity sign.
\writeup The residues modulo $11$ are $0,1,3,4,5,9$, as in the example of this chapter.
Since $13\equiv 5\pmod 8$, $\bigl(\frac{2}{13}\bigr)=-1$. Also $13\equiv 1\pmod 4$, so $\bigl(\frac{3}{13}\bigr)=\bigl(\frac{13}{3}\bigr)=\bigl(\frac{1}{3}\bigr)=1$. Hence $\bigl(\frac{6}{13}\bigr)=-1$.
Since $17\equiv 1\pmod 8$, $\bigl(\frac{2}{17}\bigr)=1$, and again $\bigl(\frac{3}{17}\bigr)=\bigl(\frac{17}{3}\bigr)=\bigl(\frac{2}{3}\bigr)=-1$, so $\bigl(\frac{6}{17}\bigr)=-1$.
Squares modulo $13$ are $0,1,3,4,9,10,12$, which do not include $6$.
\end{solution}

\begin{solution}{pb:primsum}
\idea Inside $\mathbb F_p$, a primitive root $\gamma$ plays the role of $e^{2i\pi/(p-1)}$: it has order exactly $p-1$, its
powers are all of $\mathbb F_p^*$, and the primitive roots are the $\gamma^a$ with $a$ prime to $p-1$. So the sum of their
$k$-th powers is a Ramanujan sum computed in $\mathbb F_p$, and the proof of Problem~\ref{pb:ramsum}(a) works word for word,
because it only used the geometric sum.

\writeup Let $q=p-1$ and let $\gamma$ be a primitive root (Volume~2, Chapter~4). The primitive roots modulo $p$ are the
$\gamma^a$, $1\le a\le q$, $\gcd(a,q)=1$, so in $\mathbb F_p$ the sum of their $k$-th powers is
$C_q(k)=\sum_{\gcd(a,q)=1}\gamma^{ak}$. For $e\mid q$ put $\gamma_e=\gamma^{q/e}$, an element of order $e$, and
$C_e(k)=\sum_{1\le b\le e,\ \gcd(b,e)=1}\gamma_e^{bk}$. The geometric sum $\sum_{a=1}^e\gamma_e^{ak}$ is $e$ if $e\mid k$ (all
terms are $1$) and $\frac{\gamma_e^{k}(\gamma_e^{ek}-1)}{\gamma_e^k-1}=0$ otherwise. Sorting $a$ by $\gcd(a,e)$ as in
Problem~\ref{pb:ramsum}(a), and noting that $\gamma_e^{e/f}=\gamma_f$, gives $\sum_{f\mid e}C_f(k)=e\cdot\mathbf 1_{e\mid k}$ in
$\mathbb F_p$, for all $e\mid q$. Möbius inversion over the divisors of $q$ (valid in any abelian group) gives
$C_q(k)=\sum_{d\mid\gcd(q,k)}\mu(q/d)d=c_q(k)$, which is (c).

(a) is $k=1$: $c_q(1)=\mu(q)$. For (b), $g\mapsto g^{-1}$ permutes the primitive roots, and $g\ne g^{-1}$ because $g^2=1$
would mean $q\le2$, that is $p=3$. So the primitive roots split into pairs $\{g,g^{-1}\}$ with product $1$.
\tools Volume~2, Chapters~6 and~4 (Möbius inversion, primitive roots); Volume~3, Chapter~9 (cyclic multiplicative group).
\leads Gauss stated (a) in the \emph{Disquisitiones} (article~81). Since $\mu(p-1)=0$ as soon as $4\mid p-1$, the sum of the
primitive roots is then divisible by $p$. Counting primitive roots in short intervals, where these sums meet character sums,
is still an active field (Volume~6).
\end{solution}

\begin{solution}{pb:phimodp}
\idea Modulo $p$, the polynomial $X^n-1$ still factors as $\prod_{d\mid n}\Phi_d$, and if $p\nmid n$ it has no repeated root.
So the roots of $\Phi_n$ in a big enough finite field are exactly the elements of order $n$. The degree of such an element
over $\mathbb F_p$ is the size of the smallest field $\mathbb F_{p^f}$ that contains it, and $\mathbb F_{p^f}^*$ is cyclic of order
$p^f-1$. So everything is decided by when $n$ divides $p^f-1$.

\writeup (a) Reduce $X^n-1=\prod_{d\mid n}\Phi_d$ modulo $p$. Since $p\nmid n$, the derivative $nX^{n-1}$ is prime to $X^n-1$, so
$X^n-1$ has distinct roots in a splitting field $\mathbb F$ (Volume~3, Chapter~9). A root $\zeta$ of $\Phi_n$ has order dividing
$n$; if its order were $d<n$, $\zeta$ would also be a root of $X^d-1=\prod_{e\mid d}\Phi_e$, hence of some $\Phi_e$ with $e\ne n$,
and a double root of $X^n-1$. So the $\varphi(n)$ roots of $\Phi_n$ modulo $p$ are distinct and of order exactly $n$. For such
a $\zeta$, the degree of its minimal polynomial over $\mathbb F_p$ is $[\mathbb F_p(\zeta):\mathbb F_p]$, the least $f$ with
$\zeta\in\mathbb F_{p^f}$, that is with $\zeta^{p^f-1}=1$, that is $n\mid p^f-1$: $f=\operatorname{ord}_n(p)$. So every irreducible
factor of $\Phi_n$ modulo $p$ has degree $f$, they are distinct (no repeated roots), and there are $\varphi(n)/f$ of them.

(b) $X^4+1=\Phi_8$ is irreducible over $\Q$ (Volume~3, Chapter~10, or Eisenstein after $X\mapsto X+1$). Modulo $2$,
$X^4+1=(X+1)^4$. For odd $p$, $p^2\equiv1\pmod8$, so $f=\operatorname{ord}_8(p)\le2$, and $\Phi_8$ splits into at least two factors.

(c) First let $a=1$. Since $p\nmid m$, $\Phi_m(X^p)=\Phi_{pm}(X)\Phi_m(X)$: both sides are monic of degree $p\varphi(m)$, and
the roots of the left side, the $z$ with $z^p$ of order $m$, have order $m$ or $pm$, all simple. Modulo $p$,
$\Phi_m(X^p)\equiv\Phi_m(X)^p$ (Frobenius: $(\sum c_iX^i)^p\equiv\sum c_iX^{ip}$ for $c_i\in\mathbb F_p$). Cancelling $\Phi_m$ in the
domain $\mathbb F_p[X]$, $\Phi_{pm}\equiv\Phi_m^{p-1}$. For $a\ge2$, rule (ii) of Problem~\ref{pb:phipm1} gives
$\Phi_{p^am}(X)=\Phi_{pm}(X^{p^{a-1}})\equiv\Phi_{pm}(X)^{p^{a-1}}\equiv\Phi_m^{(p-1)p^{a-1}}$.
\tools Volume~3, Chapters~9 and~10 (finite fields, cyclotomic polynomials); Volume~2, Chapter~4 (orders).
\leads (a) is the simplest case of the decomposition of primes in number fields: $p$ splits in $\Q(\zeta_n)$ into
$\varphi(n)/f$ primes of degree $f$ (Volume~13). $X^4+1$ shows that irreducibility over $\Q$ cannot always be detected modulo a
single prime, a fact that every polynomial factoring algorithm must live with.
\end{solution}

\begin{solution}{pb:fermatnum}
\idea A prime factor $p$ of $F_n$ makes $2$ have a large order, a power of $2$, which must divide $p-1$. That already gives
$2^{n+1}\mid p-1$; the extra factor $2$ comes from knowing that $2$ is a square. For $641$, the two expressions of $641$ say
that $5\cdot2^7\equiv-1$ and $2^4\equiv-5^4$, and multiplying the right powers gives $2^{32}\equiv-1$.

\writeup (a) Let $p\mid F_n$. Then $2^{2^n}\equiv-1$, so $2^{2^{n+1}}\equiv1$ and the order of $2$ modulo $p$ divides $2^{n+1}$ but
not $2^n$: it is $2^{n+1}$, and it divides $p-1$ (Volume~2, Chapter~4). For $n\ge2$ this gives $p\equiv1\pmod8$, so
$\bigl(\frac2p\bigr)=1$ (Volume~2, Chapter~5) and, by Euler's criterion, $2^{(p-1)/2}\equiv1$. So $2^{n+1}\mid\frac{p-1}2$, and
$p\equiv1\pmod{2^{n+2}}$.

(b) Modulo $641$: $5\cdot2^7\equiv-1$, so $5^4\cdot2^{28}\equiv1$; and $2^4\equiv-5^4$. Hence
$2^{32}=2^4\cdot2^{28}\equiv-5^4\cdot2^{28}\equiv-1$, and $641\mid F_5$. In accordance with (a), $641=10\cdot2^6+1\equiv1\pmod{2^7}$: Euler
needed only a few trials.

(c) $F_0F_1\cdots F_{n-1}=F_n-2$, by induction: $(F_n-2)F_n=(2^{2^n}-1)(2^{2^n}+1)=2^{2^{n+1}}-1=F_{n+1}-2$. So a common divisor
of $F_m$ and $F_n$, $m<n$, divides $2$; the $F_n$ are odd, so they are coprime. Each $F_n$ has a prime factor, and these are
all different: there are infinitely many primes.
\tools Volume~2, Chapters~4 and~5 (orders, the value of $(\frac2p)$, Euler's criterion).
\leads Fermat conjectured in 1640 that all $F_n$ are prime; Euler found $641$ in 1732. Only $F_0,\dots,F_4$ are known to be prime,
and $F_5,\dots,F_{32}$ are known to be composite. Landry found $F_6=274177\cdot67280421310721$ in 1880, and indeed
$274177\equiv1\pmod{2^8}$. Gauss proved that the regular $n$-gon is constructible when $n$ is a power of $2$ times distinct Fermat
primes (Volume~3, Chapter~10).
\end{solution}

\begin{solution}{pb:carmichael}
\idea To show that $n$ fails, find one base $a$ that breaks $a^{n-1}\equiv1$. The Chinese remainder theorem lets us choose
$a$ freely modulo each prime power factor, so we can take it of maximal order there. Then $a^{n-1}\equiv1$ forces that order
to divide $n-1$. With two prime factors, the divisibilities $p-1\mid n-1$ and $q-1\mid n-1$ cannot both hold.

\writeup (a) $n$ is odd: otherwise $(-1)^{n-1}=-1\not\equiv1\pmod n$, as $n>2$. Let $p\mid n$ and write $n=p^vm$ with $p\nmid m$.
If $v\ge2$, let $g$ have order $p(p-1)$ modulo $p^2$ ($(\Z/p^2\Z)^*$ is cyclic, Volume~2, Chapter~4), and by the Chinese
remainder theorem choose $a\equiv g\pmod{p^2}$, $a\equiv1\pmod m$, and $a$ prime to $p$; then $a$ is prime to $n$. From
$a^{n-1}\equiv1\pmod{p^2}$, $p(p-1)\mid n-1$, so $p\mid n-1$, absurd since $p\mid n$. So $n$ is squarefree. Now take $a\equiv g$
modulo $p$ with $g$ a primitive root, and $a\equiv1\pmod m$: $a^{n-1}\equiv1\pmod p$ gives $p-1\mid n-1$.

(b) $n$ is squarefree and composite, so it has at least two prime factors. If $n=pq$ with $p<q$, then
$n-1=p(q-1)+(p-1)\equiv p-1\pmod{q-1}$, and $q-1\mid n-1$ would give $q-1\mid p-1$, impossible as $0<p-1<q-1$.

(c) $n=(6k+1)(12k+1)(18k+1)=1296k^3+396k^2+36k+1$, so $n-1=36k(36k^2+11k+1)$ is divisible by $6k$, $12k$ and $18k$. By
Korselt's criterion (Volume~2, Chapter~3), $n$ is a Carmichael number. For $k=1$: $7\cdot13\cdot19=1729$.
\tools Volume~2, Chapters~3 and~4 (Chinese remainder theorem, Korselt's criterion, primitive roots modulo $p^2$).
\leads $1729$ is Hardy's taxi-cab number, Ramanujan's ``smallest number expressible as a sum of two cubes in two different
ways''. Chernick (1939) gave (c); whether it gives infinitely many examples is open, but Alford, Granville and Pomerance proved
in 1994 that there are infinitely many Carmichael numbers.
\end{solution}

\begin{solution}{pb:pratt}
\idea (a) An element whose order is exactly $n-1$ forces the group $(\Z/n\Z)^*$ to have $n-1$ elements, which only happens
for primes; and checking that the order is $n-1$ needs only the prime factors of $n-1$. (b) The tree branches at the prime
factors of $p-1$, whose product is less than $p$: induction on $p$ with a logarithmic bound, because logarithms add.

\writeup (a) If $n$ is prime, a primitive root works (Volume~2, Chapter~4). Conversely, the order $e$ of $a$ in $(\Z/n\Z)^*$ ($a$ is a
unit since $a\cdot a^{n-2}\equiv1$) divides $n-1$. If $e<n-1$, some prime $q$ divides $\frac{n-1}e$, and then $e\mid\frac{n-1}q$ and
$a^{(n-1)/q}\equiv1$, excluded. So $e=n-1$, hence $n-1\le\varphi(n)$, and $n$ is prime.

(b) Induction on $p$. As $p$ is odd, $2\mid p-1$. Let $q_1,\dots,q_k$ be the distinct odd primes dividing $p-1$; their product is at
most $\frac{p-1}2$. If $k=0$, $f(p)=1+f(2)=2\le2\log_2p-1$, because $p\ge3$. If $k\ge1$, by induction
\[
f(p)=1+f(2)+\sum_{i=1}^kf(q_i)\le2+\sum_i(2\log_2q_i-1)\le2+2\log_2\frac{p-1}2-k=2\log_2(p-1)-k,
\]
which is at most $2\log_2p-1$.

(c) A prime $p$ with $k$ digits has at most $2k$ primes in its tree, each with at most $k$ digits. At each one, say $\ell$, the
checker multiplies the claimed factorisation of $\ell-1$ (at most $k$ factors with multiplicity) and computes $a^{\ell-1}$ and the
$a^{(\ell-1)/q}$ modulo $\ell$, at most $k+1$ modular powers, each by $O(k)$ squarings and multiplications (Volume~2, Chapter~3). In
total $O(k^3)$ operations on numbers of $k$ digits.
\tools Volume~2, Chapters~3 and~4 (fast exponentiation, orders, primitive roots).
\leads Pratt (1975): \emph{PRIMES is in NP}, the first complexity result on primality. Finding the certificate requires factoring
$p-1$, which is hard in general; but primes of special shape (Proth primes, Problem~\ref{pb:pocklington}) carry their own
factorisation, and elliptic curve methods replace $p-1$ by the order of a random group (Goldwasser--Kilian, Atkin--Morain).
\end{solution}

\begin{solution}{pb:rabin}
\idea A strong liar is described componentwise by the Chinese remainder theorem, and in each cyclic group $(\Z/p_i^{e_i}\Z)^*$
the equations $x^t=\pm1$ have explicit numbers of solutions. So the liars can be counted exactly, and the count is then
compared with $\varphi(n)$ factor by factor. The prime powers lose a factor $p^{e-1}$; several primes lose a factor
$2^{r-1}$ from the powers of $2$; and the only delicate case, two primes with the same $2$-part, is saved by a
divisibility argument.

\writeup (a) Let $g$ generate $C$, of order $m$. $x=g^u$ satisfies $x^t=1$ iff $m\mid tu$, iff $\frac m{\gcd(t,m)}\mid u$: there are
$\gcd(t,m)$ such $u$ modulo $m$. If $x^t=-1$, then $x^{2t}=1$ and $x^t\ne1$. Conversely if $x^{2t}=1\ne x^t$, then $x^t$ has order $2$,
so it is $-1$, the unique element of order $2$ in a cyclic group of even order. So the count is $\gcd(2t,m)-\gcd(t,m)$.

(b) Let $\varphi_i=p_i^{e_i-1}2^{\nu_i}m_i$, the order of the cyclic group $C_i=(\Z/p_i^{e_i}\Z)^*$ (Volume~2, Chapter~4; $p_i$ is odd).
By the Chinese remainder theorem, $a$ is a liar iff either $a^d=1$ in every $C_i$, or for some $0\le j<s$, $a^{2^jd}=-1$ in every
$C_i$. These cases are disjoint: if $a^{2^jd}\equiv-1$, then $a^{2^{j'}d}\equiv1$ for $j'>j$, and $\pm1$ are distinct modulo $n$. Since
$p_i\mid n$, $p_i\nmid n-1$, so $p_i\nmid d$, and $\gcd(2^jd,\varphi_i)=g_i\,2^{\min(j,\nu_i)}$ with $g_i=\gcd(d,m_i)$ ($d$ is odd). By (a):
\begin{itemize}
\item $a^d=1$ in $C_i$ for $g_i$ elements;
\item $a^{2^jd}=-1$ in $C_i$ for $g_i\bigl(2^{\min(j+1,\nu_i)}-2^{\min(j,\nu_i)}\bigr)$ elements, which is $g_i2^j$ if $j<\nu_i$ and $0$ if
$j\ge\nu_i$.
\end{itemize}
Multiplying over $i$, the number of $a$ with $a^{2^jd}\equiv-1\pmod n$ is $\bigl(\prod_ig_i\bigr)2^{jr}$ if $j<\nu$ and $0$ otherwise.
Each $p_i\equiv1\pmod{2^\nu}$, so $n\equiv1\pmod{2^\nu}$ and $\nu\le s$. Summing over $0\le j<\nu$ and adding the first case, the
number of liars is $\prod_ig_i\cdot\bigl(1+\sum_{j<\nu}2^{jr}\bigr)$, which is Monier's formula.

(c) Write $T_r(\nu)=\frac{1+\frac{2^{r\nu}-1}{2^r-1}}{2^{r\nu}}=2^{-r\nu}\Bigl(1-\frac1{2^r-1}\Bigr)+\frac1{2^r-1}$. Then the proportion of liars is
\[
R=\frac{L}{\varphi(n)}=\prod_{i=1}^r\frac{g_i}{m_ip_i^{e_i-1}}\cdot T_r(\nu)\cdot2^{r\nu-(\nu_1+\dots+\nu_r)}.
\]
All three factors are at most $1$, and $T_r$ decreases in $\nu\ge1$, so $T_r(\nu)\le T_r(1)=2^{-r}+2^{-r}=2^{1-r}$ (for $r=1$,
$T_1(\nu)=1$).

\emph{$r=1$}, $n=p^e$ with $e\ge2$: $R\le p^{1-e}$, which is $\frac13$ for $n=9$ and at most $\frac15$ otherwise ($p^{e-1}\ge5$).

\emph{$r\ge3$}: $R\le2^{1-r}\le\frac14$.

\emph{$r=2$}: $R\le\frac12\cdot\frac{g_1g_2}{m_1m_2p_1^{e_1-1}p_2^{e_2-1}}\cdot2^{2\nu-\nu_1-\nu_2}$. If $\nu_1\ne\nu_2$, the last factor is at most
$\frac12$. If some $e_i\ge2$, the middle factor is at most $\frac13$. There remains $n=p_1p_2$ with $\nu_1=\nu_2=\nu$. If
$g_1=m_1$ and $g_2=m_2$, then $m_i\mid d$, so $p_i-1=2^\nu m_i$ divides $2^sd=n-1$. But $n-1=p_1(p_2-1)+(p_1-1)$, so $p_2-1\mid p_1-1$,
and symmetrically $p_1-1\mid p_2-1$, so $p_1=p_2$: absurd. So some $g_i$ is a proper divisor of the odd $m_i$, $\frac{g_i}{m_i}\le\frac13$,
and $R\le\frac16$.

In all cases $R\le\frac14$, except $n=9$. For $n=15$: $14=2\cdot7$, $s=1$, $d=7$; $r=2$, $p_1-1=2$, $p_2-1=4$, $\nu=1$, $g_1=g_2=1$. Monier
gives $1\cdot(1+1)=2$ liars ($\pm1$) out of $\varphi(15)=8$: exactly $\frac14$.
\tools Volume~2, Chapters~3 and~4 (Chinese remainder theorem, cyclic groups $(\Z/p^e\Z)^*$); Volume~3, Chapter~1 (cyclic groups).
\leads Miller (1976) proved that under the generalised Riemann hypothesis the bases $a<2\log^2n$ suffice, a deterministic test;
Rabin (1980) turned it into the randomised test with error $4^{-k}$ for $k$ bases. Figure~\ref{fig:liars} shows that the
bound is far from typical: for most $n$ the proportion of liars is tiny, which is why, in practice, a handful of bases is
used for numbers of thousands of digits.
\end{solution}

\begin{solution}{pb:waring}
\idea The number of $k$-th roots of $u$ is a sum over the $k$ characters of order dividing $k$. Then $N(c)$ is $p$ plus a sum of
$(k-1)^2$ Jacobi sums, of which $k-1$ have size $1$ and the others size $\sqrt p$.

\writeup (a) $\mathbb F_p^*$ is cyclic and $k\mid p-1$, so the characters $\chi$ with $\chi^k=\varepsilon$ form a group of order $k$,
and for $u\ne0$, $\sum_{\chi^k=\varepsilon}\chi(u)$ is $k$ if $u$ is a $k$-th power and $0$ otherwise: the number of $x$ with $x^k=u$.
With $\varepsilon(0)=1$ and $\chi(0)=0$ for $\chi\ne\varepsilon$ it is also right for $u=0$. So
\[
N(c)=\sum_{\chi,\psi}\sum_{u+v=c}\chi(u)\psi(v)=\sum_{\chi,\psi}\chi\psi(c)J(\chi,\psi),
\]
after the substitution $u=cu'$, $v=cv'$. The pair $(\varepsilon,\varepsilon)$ gives $p$; pairs with exactly one trivial character
give $0$; the $k-1$ pairs $(\chi,\bar\chi)$ with $\chi\ne\varepsilon$ give $J(\chi,\bar\chi)=-\chi(-1)$, of absolute value $1$ (as in
Problem~\ref{pb:fermatcubic}(a)); and the remaining $(k-1)^2-(k-1)=(k-1)(k-2)$ pairs have $\chi,\psi,\chi\psi$ all non-trivial,
so $|J(\chi,\psi)|=|g(\chi)||g(\psi)|/|g(\chi\psi)|=\sqrt p$ by Problem~\ref{pb:fermatcubic}(b), whose proof works for any characters.

(b) $c=0$ is $0^k+0^k$. For $c\ne0$, $N(c)>0$ if $p>(k^2-3k+2)\sqrt p+k-1$. If $p\ge k^4$, then $\sqrt p\ge k^2$ and
$p-(k^2-3k+2)\sqrt p-(k-1)\ge(3k-2)\sqrt p-(k-1)>0$.

(c) If $d=\gcd(k,p-1)$, the $k$-th powers of $\mathbb F_p^*$ are the $d$-th powers (the image of $x\mapsto x^k$ in the cyclic group
of order $p-1$ is the subgroup of index $d$). So the question reduces to $d$; when $d=1$, every residue is a $k$-th power.

\tools Volume~2, Chapter~4 (characters of $\mathbb F_p^*$); Problem~\ref{pb:fermatcubic}.
\leads For $k=3$ the bound is $2\sqrt p+2$, matching Problem~\ref{pb:fermatcubic}. Weil's theorem bounds the error by
$(k-1)(k-2)\sqrt p$ for the projective curve, of genus $\frac{(k-1)(k-2)}2$. A numerical calculation lists the primes where two $k$-th powers
miss a residue: for $k=3$ only $p=7$, for $k=4$ the largest is $29$, for $k=8$ it is $761$. They are all well below $k^4$, so the
exponent $4$ is far from sharp.
\end{solution}

\begin{solution}{pb:squaresmodpk}
\idea A non-zero square modulo $p^k$ has an even valuation $2j<k$ and a unit part that is a square modulo $p^{k-2j}$. Counting
units that are squares is a matter of group theory: the squaring map on the cyclic unit group has kernel $\pm1$.

\writeup (a) Let $p$ be odd. If $x=y^2\not\equiv0\pmod{p^k}$, then $v_p(x)=2v_p(y)=2j<k$ and $x=p^{2j}u$ with $u$ a unit modulo $p^{k-2j}$ which is a
square modulo $p^{k-2j}$. Conversely every such $p^{2j}u$ is a square. The residue $x$ modulo $p^k$ determines $j$ and $u$ modulo $p^{k-2j}$,
and the units modulo $p^m$ form a cyclic group of even order $\varphi(p^m)$ (Volume~2, Chapter~4), in which the squares are a subgroup of
index $2$. So there are $\frac12\varphi(p^{k-2j})=\frac{p-1}2p^{k-2j-1}$ non-zero squares of valuation $2j$, and adding $0$ gives the formula.
Dividing by $p^k$, $\frac{S(p^k)}{p^k}=p^{-k}+\frac{p-1}{2p}\sum_{0\le j<k/2}p^{-2j}\to\frac{p-1}{2p}\cdot\frac1{1-p^{-2}}=\frac p{2(p+1)}$.

(b) For $p=2$ the units modulo $2^m$ that are squares are those $\equiv1\pmod8$ when $m\ge3$ (and only $1$ for $m\le2$): a proportion
$\frac14$ of the $2^{m-1}$ units, that is $2^{m-3}$. So $\frac{S(2^k)}{2^k}=\sum_{2j\le k-3}2^{-2j}\cdot\frac18+O(2^{-k/2})\to\frac18\cdot\frac43=\frac16$.

(c) By the Chinese remainder theorem, $x$ is a square modulo $mn$ (coprime) if and only if it is a square modulo $m$ and modulo $n$,
so $S(mn)=S(m)S(n)$. For $m=\prod_{p\le y}p$, the proportion of squares is $\prod_{p\le y}\frac{p+1}{2p}\le\prod_{3\le p\le y}\frac23\to0$. An
integer $a$ that is a square modulo every $m$ is a square: otherwise write $a=b^2c$ with $c\ne1$ squarefree (or $a<0$), and choose,
by Chinese remainders and quadratic reciprocity, a prime $q\nmid2a$ with $\bigl(\frac cq\bigr)=-1$; then $a$ is not a square modulo $q$.

\tools Volume~2, Chapters~4 and~5, and Volume~13, Chapter~8 (cyclic unit groups, reciprocity, $p$-adic squares); Volume~3, Chapter~2.
\leads In $p$-adic language, the squares of $\Z_p$ have Haar measure $\frac p{2(p+1)}$ ($\frac16$ for $p=2$). The last statement of (c)
is the simplest case of a \emph{local--global principle}: a property that holds modulo every $m$ is forced
to hold over $\Z$. Such principles are a recurring theme of Part~III.
\end{solution}

\section*{Sources and further reading}

\begin{itemize}
\item K.~Ireland and M.~Rosen, \emph{A Classical Introduction to Modern Number Theory}, 2nd ed., Springer, 1990, Chapters~5--6. Reciprocity, Gauss sums and the beginning of cyclotomy.
\item F.~Lemmermeyer, \emph{Reciprocity Laws: From Euler to Eisenstein}, Springer, 2000. The full history, with a list of all known proofs of quadratic reciprocity (kept up to date on the author's web page).
\item D.~A.~Cox, \emph{Primes of the Form $x^2+ny^2$}, 3rd ed., AMS Chelsea, 2022, \S1. How reciprocity answers Fermat's questions, the theme of Chapters~7 and~10.
\item R.~Solovay and V.~Strassen, A fast Monte-Carlo test for primality, \emph{SIAM Journal on Computing} 6 (1977), 84--85. The Jacobi symbol at work in the first randomised primality test.
\item A.~Granville, It is easy to determine whether a given integer is prime, \emph{Bulletin of the American Mathematical Society} 42 (2005), 3--38. A survey of where primality testing went next, up to the AKS algorithm.
\item J.~J.~Gray, \emph{A commentary on Gauss's mathematical diary}, Expositiones Mathematicae 2 (1984), 97--130, entry of 8 April 1796; F.~Lemmermeyer, \emph{Reciprocity Laws}, Springer, 2000, whose appendix lists the published proofs.
\end{itemize}
